\documentclass[aspectratio=169,11pt]{beamer}

\usepackage{iftex}
\ifPDFTeX
  \usepackage[utf8]{inputenc}
  \usepackage[T5]{fontenc}
  \usepackage[vietnamese]{babel}
\else
  \usepackage{fontspec}
  \setsansfont{TeX Gyre Heros}
  \setmainfont{TeX Gyre Termes}
  \setmonofont{Inconsolata}
  \usepackage[vietnamese]{babel}
\fi

\usepackage{graphicx}
\usepackage{booktabs}
\usepackage{tikz}
\usepackage{listings}
\usepackage{xcolor}
\usepackage{hyperref}
\usetikzlibrary{shapes.geometric, arrows.meta, positioning, shadows}
\usepackage{pifont}

\vfuzz=120pt
\hbadness=10000
\hfuzz=2pt
\pdfstringdefDisableCommands{%
  \def\textbf#1{#1}%
  \def\\{ }%
}
\DeclareFontShape{T5}{cmss}{m}{it}{<->ssub*cmss/m/sl}{}

\definecolor{hcmusblue}{HTML}{005BAA}
\definecolor{hcmusred}{HTML}{D71920}
\definecolor{darktext}{HTML}{202428}
\definecolor{primaryblue}{RGB}{41, 128, 185}
\definecolor{softgray}{RGB}{236, 240, 241}
\definecolor{accentorange}{RGB}{230, 126, 34}
\definecolor{successgreen}{RGB}{39, 174, 96}
\definecolor{dangerred}{RGB}{192, 57, 43}

\usetheme{Madrid}
\usecolortheme{default}
\setbeamercolor{structure}{fg=hcmusblue}
\setbeamercolor{title}{fg=white,bg=hcmusblue}
\setbeamercolor{frametitle}{fg=white,bg=hcmusblue}
\setbeamercolor{block title}{fg=white,bg=hcmusblue}
\setbeamercolor{block body}{bg=softgray,fg=darktext}
\setbeamertemplate{navigation symbols}{}
\setbeamertemplate{footline}[frame number]

\newcommand{\ProjectTitle}{Kiến trúc hệ RAG phát hiện đạo văn}
\newcommand{\CourseName}{Khai thác dữ liệu văn bản và ứng dụng}
\newcommand{\HCMUSLogo}{%
  \IfFileExists{figures/hcmus-logo.png}{%
    \includegraphics[height=0.9cm]{figures/hcmus-logo.png}%
  }{%
    \begin{tikzpicture}[scale=0.36]
      \draw[line width=1.2pt,color=hcmusblue] (0,0) circle (1.25);
      \draw[line width=0.8pt,color=hcmusred] (0,0) circle (0.95);
      \node[font=\bfseries\scriptsize,color=hcmusblue] at (0,0) {HCMUS};
    \end{tikzpicture}%
  }%
}

% ---- TikZ styles dùng chung ----
\tikzset{
  box/.style={draw=primaryblue!80, fill=primaryblue!10, thick, rounded corners=4pt,
    align=center, minimum height=0.85cm, font=\scriptsize\bfseries},
  gbox/.style={draw=successgreen!80, fill=successgreen!10, thick, rounded corners=4pt,
    align=center, minimum height=0.85cm, font=\scriptsize\bfseries},
  obox/.style={draw=accentorange!85, fill=accentorange!12, thick, rounded corners=4pt,
    align=center, minimum height=0.85cm, font=\scriptsize\bfseries},
  nbox/.style={draw=gray!55, fill=gray!8, thick, rounded corners=4pt,
    align=center, minimum height=0.7cm, font=\scriptsize},
  ar/.style={-{Latex[scale=1.1]}, line width=1pt, primaryblue},
  arO/.style={-{Latex[scale=1.1]}, line width=1.3pt, accentorange},
  slbl/.style={font=\tiny, text=black!70},
}

\title[\ProjectTitle]{\textbf{\ProjectTitle}}
\subtitle{\CourseName}

\author[Nhóm 3]{\texorpdfstring{%
  \small\textbf{GVHD: Ph.D Lê Thanh Tùng}\\[0.25cm]
  \textbf{Nhóm 3}\\[0.1cm]
  \scriptsize
  \begin{tabular}{l}
    \textbf{Sinh viên thực hiện:} \\
    23127004 - Lê Nhật Khôi \\
    23127049 - Đỗ Hoàng Duy Hưng \\
    23127189 - Trần Trọng Hiếu \\
    23127210 - Nguyễn Đăng Khoa \\
    23127248 - Nguyễn Hữu Phúc
  \end{tabular}
}{Nhóm 3}}
\hypersetup{pdfauthor={Nhóm 3},pdftitle={Kiến trúc hệ RAG phát hiện đạo văn}}

\institute[HCMUS]{
  Khoa Công nghệ thông tin\\
  Trường Đại học Khoa học tự nhiên, ĐHQG-HCM
}
\date{Tháng 08 năm 2026}

\begin{document}

\begin{frame}
  \titlepage
  \begin{tikzpicture}[remember picture,overlay]
    \node[anchor=north east,xshift=-0.55cm,yshift=-0.85cm] at (current page.north east) {\HCMUSLogo};
  \end{tikzpicture}
\end{frame}

\begin{frame}{Nội dung}
\tableofcontents
\end{frame}

% ============================================================
\section{Kiến trúc pipeline}
{
\setbeamercolor{background canvas}{bg=white}
\begin{frame}[plain]
    \vfill\centering
    \textcolor{hcmusblue}{\textbf{\Huge Kiến trúc pipeline}}
    \vfill
\end{frame}
}

\begin{frame}{Tổng quan — pipeline RAG năm khâu}
\centering
\begin{tikzpicture}[node distance=0.55cm and 0.7cm]
  \node[box, text width=1.7cm] (doc) {Tài liệu\\nghi vấn};
  \node[box, text width=1.9cm, right=of doc, draw=primaryblue] (ret) {Retrieval\\{\tiny TF-IDF top-k}};
  \node[obox, text width=1.9cm, right=of ret] (ali) {Alignment\\{\tiny seed-extend}};
  \node[box, text width=1.8cm, right=of ali] (ver) {Verifier\\{\tiny lọc báo giả}};
  \node[gbox, text width=2.0cm, right=of ver] (gen) {Generation\\{\tiny luận giải RAG}};
  \node[nbox, text width=1.9cm, below=0.7cm of ret] (corp) {Kho nguồn\\7\,949 tài liệu};
  \draw[ar] (doc) -- (ret);
  \draw[ar] (ret) -- node[slbl,above]{top-k} (ali);
  \draw[ar] (ali) -- node[slbl,above]{span} (ver);
  \draw[ar] (ver) -- node[slbl,above]{GIỮ} (gen);
  \draw[ar] (corp) -- (ret);
\end{tikzpicture}

\vspace{0.35cm}
\setbeamercolor{block title}{bg=primaryblue,fg=white}
\begin{block}{Bài toán PAN 2025 — Generative Plagiarism Detection}
\small
Với mỗi \textbf{tài liệu nghi vấn} và một \textbf{kho nguồn}: tìm mọi đoạn đạo văn, chỉ ra nguồn
tương ứng, xác minh, và sinh \textbf{luận giải grounded}. Chấm điểm ở \textbf{mức ký tự} bằng \texttt{PlagDet}.
Mỗi khâu tách bạch để \textbf{đánh giá độc lập}.
\end{block}
\end{frame}

\begin{frame}{Bài toán, Input -- Output \& nền dữ liệu}
\setbeamercolor{block title}{bg=primaryblue,fg=white}
\begin{block}{Input -- Output}
\begin{columns}[T]
\begin{column}{0.48\textwidth}
\textbf{Input}
\begin{itemize}\small
  \item Tài liệu nghi vấn $P$ (\texttt{.txt}).
  \item Kho nguồn $\mathcal{S}$ (7\,949 tài liệu val).
\end{itemize}
\end{column}
\begin{column}{0.48\textwidth}
\textbf{Output}
\begin{itemize}\small
  \item Span $(\text{offset}, \text{length}, \text{nguồn})$.
  \item Luận giải grounded (kỹ thuật · giải thích · viết lại).
\end{itemize}
\end{column}
\end{columns}
\end{block}

\vspace{0.1cm}
\setbeamercolor{block title}{bg=accentorange,fg=white}
\begin{block}{Các tầng của hệ thống}
\centering\scriptsize
\renewcommand{\arraystretch}{1.15}
\begin{tabular}{llll}
\toprule
\textbf{Tầng} & \textbf{Đầu vào} & \textbf{Đầu ra} & \textbf{Phương pháp} \\
\midrule
Retrieval  & tài liệu nghi vấn & top-k nguồn khả nghi & TF-IDF + cosine \\
Alignment  & cặp (nghi vấn, nguồn) & span đạo + span nguồn & tf-isf seed-and-extend \\
Scoring    & span vs gold & PlagDet / P / R & PAN 2015 char-level \\
Verifier   & span đã gắn cờ & GIỮ / BỎ (lọc FP) & LLM verifier (GLM-5.2) \\
Generation & span + nguồn truy hồi & luận giải grounded & GLM-5.2 (retrieval-augmented) \\
\bottomrule
\end{tabular}
\end{block}
\centering\footnotesize\textbf{Dữ liệu:} 7\,949 nguồn · 5\,522 nghi vấn · arXiv (ar5iv) · thước đo \texttt{PlagDet} (PAN 2015).
\end{frame}

% ============================================================
\section{Retrieval}
{
\setbeamercolor{background canvas}{bg=white}
\begin{frame}[plain]
    \vfill\centering
    \textcolor{hcmusblue}{\textbf{\Huge Không gian tìm kiếm \& Retrieval}}
    \vfill
\end{frame}
}

\begin{frame}{Retrieval — thu hẹp 7\,949 nguồn bằng TF-IDF}
\begin{columns}[T]
\begin{column}{0.55\textwidth}
\setbeamercolor{block title}{bg=primaryblue,fg=white}
\begin{block}{Xây chỉ mục}
\scriptsize
\begin{enumerate}\setlength\itemsep{1pt}
  \item \textbf{Chunking theo ký tự:} \texttt{20\,000} ký tự đầu mỗi tài liệu $\rightarrow$ 1 vector/doc.
  \item \textbf{Vector hoá TF-IDF:} \texttt{max\_features=100k}, \texttt{sublinear\_tf}, bỏ stop-word EN.
  \item \textbf{Truy vấn:} nghi vấn $\rightarrow$ vector $\rightarrow$ \textbf{cosine} toàn kho $\rightarrow$ top-k.
\end{enumerate}
\vspace{2pt}\scriptsize \textbf{In:} tài liệu + chỉ mục TF-IDF $\;\rightarrow\;$ \textbf{Out:} top-k nguồn xếp theo cosine.
\end{block}
\end{column}
\begin{column}{0.43\textwidth}
\centering
\begin{tikzpicture}[baseline]
  \node[box, text width=1.4cm] (q) at (0,0) {vector\\nghi vấn};
  \node[font=\tiny, text=gray!80] at (3.0,0.95){kho 7\,949 vector nguồn};
  \foreach \y in {0.62,0.40,0.18,-0.04,-0.26} \draw[fill=gray!12,draw=gray!45] (2.2,\y) rectangle (3.8,\y+0.16);
  \draw[fill=accentorange,draw=accentorange] (2.2,-0.62) rectangle (3.8,-0.44);
  \node[font=\tiny,text=white] at (3.0,-0.53){\#1 nguồn khớp nhất};
  \draw[ar] (q.east) -- node[slbl,above]{cosine} (2.15,0.1);
\end{tikzpicture}
\\[0.15cm]
{\scriptsize Retrieval $=$ nhân ma trận thưa $\rightarrow$ xếp hạng cosine.}
\end{column}
\end{columns}

\vspace{0.05cm}
\setbeamercolor{block title}{bg=successgreen,fg=white}
\begin{block}{Kết quả retrieval — val 5\,522 susp (\textbf{TF-IDF thắng neural})}
\centering\scriptsize
\renewcommand{\arraystretch}{1.05}
\begin{tabular}{lrrrr}
\toprule
\textbf{Phương pháp} & \textbf{R@1} & \textbf{R@5} & \textbf{R@10} & \textbf{MRR} \\
\midrule
\textbf{TF-IDF} & \textbf{0.975} & \textbf{0.994} & \textbf{0.995} & \textbf{0.985} \\
Neural (bge-small) & 0.908 & 0.965 & 0.975 & 0.934 \\
\bottomrule
\end{tabular}
\end{block}
\centering\footnotesize Đạo văn PAN trùng lặp từ vựng cao — lexical bắt trực tiếp cụm từ dùng chung; embedding làm mờ tài liệu cùng chủ đề.
\end{frame}

% ============================================================
\section{Alignment}
{
\setbeamercolor{background canvas}{bg=white}
\begin{frame}[plain]
    \vfill\centering
    \textcolor{hcmusblue}{\textbf{\Huge Alignment — đối chiếu từng đoạn}}
    \vfill
\end{frame}
}

\begin{frame}{Alignment — seed-and-extend (Sánchez-Pérez, PAN 2014)}
\begin{columns}[T]
\begin{column}{0.56\textwidth}
\setbeamercolor{block title}{bg=primaryblue,fg=white}
\begin{block}{Cơ chế}
\scriptsize
\begin{enumerate}\setlength\itemsep{1pt}
  \item \textbf{Cắt câu giữ offset:} mỗi câu biết vị trí ký tự.
  \item \textbf{Vector tf-isf + Dice:} ma trận cosine \& Dice mọi cặp câu (nghi vấn $\times$ nguồn).
  \item \textbf{Gieo hạt (seed):} hạt khi $\cos \ge \theta_1$ \textbf{và} $\mathrm{Dice} \ge \theta_2$.
  \item \textbf{Cụm hai chiều:} gộp hạt liền mạch trên đường chéo $\rightarrow$ "case".
  \item \textbf{Mở rộng \& lọc:} giữ case $\ge \theta_3$, dài $\ge 150$ ký tự; khử chồng lấn.
\end{enumerate}
\end{block}
\end{column}
\begin{column}{0.42\textwidth}
\centering
\begin{tikzpicture}[scale=0.68, every node/.style={font=\tiny}]
  \draw[gray!40] (0,0) grid (5,5);
  \node at (2.5,5.4){câu nguồn $\rightarrow$};
  \node[rotate=90] at (-0.5,2.5){câu nghi vấn $\rightarrow$};
  \fill[primaryblue!45] (0.5,4.5) circle (0.12);
  \fill[primaryblue!45] (4.5,0.5) circle (0.12);
  \foreach \x in {1,2,3} \fill[primaryblue] (\x+0.5,4.5-\x) circle (0.16);
  \fill[primaryblue] (4.5,0.5) circle (0.16);
  \draw[accentorange, line width=1.5pt, rounded corners=3pt] (1.1,1.1) rectangle (4.9,3.9);
  \node[text=accentorange] at (2.5,-0.55){hạt liền mạch $=$ 1 đoạn đạo};
\end{tikzpicture}
\\[0.15cm]{\scriptsize \textcolor{primaryblue}{$\bullet$} hạt · \textcolor{accentorange}{$\Box$} case (span).\\ In: cặp (nghi vấn, 1 nguồn) $\rightarrow$ Out: (span nghi vấn, span nguồn).}
\end{column}
\end{columns}
\end{frame}

\begin{frame}{Alignment — kết quả detection (PlagDet trên val)}
\setbeamercolor{block title}{bg=successgreen,fg=white}
\begin{block}{Cấu hình \& PlagDet}
\centering\small
\renewcommand{\arraystretch}{1.1}
\begin{tabular}{lrl}
\toprule
\textbf{Cấu hình} & \textbf{PlagDet} & \textbf{Ghi chú} \\
\midrule
\textbf{tf-isf top-1 (chính thức)} & \textbf{0.713} & full 5\,522 · 83\% của trần câu 0.858 \\
baseline (gán cả tài liệu) & 0.675 & full 5\,522 \\
tf-isf top-3 & 0.620 & distractor phá precision \\
E5+path (neural, th 0.78) & 0.663 & subset 1\,000 (giới hạn GPU) \\
E5+path (neural, th 0.86) & 0.463 & subset 1\,000 (giới hạn GPU) \\
\bottomrule
\end{tabular}
\end{block}

\vspace{0.1cm}
\setbeamercolor{block title}{bg=dangerred,fg=white}
\begin{block}{Ba phát hiện}
\footnotesize
\begin{itemize}\setlength\itemsep{1pt}
  \item[\ding{204}] \textbf{top-1 $>$ top-3} (0.713 vs 0.620): thêm nguồn distractor phá precision.
  \item[\ding{205}] \textbf{lexical $>$ neural}: tf-isf 0.713 vượt xa E5+path — trùng lặp từ vựng, embedding thua.
  \item[\ding{206}] Hạ ngưỡng E5 (0.86$\rightarrow$0.78) cứu recall (0.463$\rightarrow$0.663) nhưng chưa chạm baseline lexical.
\end{itemize}
\end{block}
\end{frame}

% ============================================================
\section{Verifier LLM}
{
\setbeamercolor{background canvas}{bg=white}
\begin{frame}[plain]
    \vfill\centering
    \textcolor{hcmusblue}{\textbf{\Huge Verifier LLM — khử báo giả}}
    \vfill
\end{frame}
}

\begin{frame}{Verifier LLM — so sánh 6 model (FPT)}
\setbeamercolor{block title}{bg=primaryblue,fg=white}
\begin{block}{Bài toán}
\footnotesize Aligner đôi khi gắn cờ hai đoạn chỉ \textbf{trùng chủ đề}. LLM đọc lại cặp (nghi vấn, nguồn),
quyết định \textbf{GIỮ / BỎ} $\rightarrow$ tăng precision. Nhãn TP/FP suy từ \textbf{gold PAN} (khách quan).
\end{block}

\vspace{0.05cm}
\setbeamercolor{block title}{bg=accentorange,fg=white}
\begin{block}{1\,000 span val · 824 TP / 176 FP · PlagDet nền 0.715}
\centering\scriptsize
\renewcommand{\arraystretch}{1.12}
\begin{tabular}{lrrrrrr}
\toprule
\textbf{Model} & \textbf{PlagDet} & \textbf{$\Delta$} & \textbf{fp\_red} & \textbf{tp\_ret} & \textbf{prec} & \textbf{lỗi} \\
\midrule
\textbf{GLM-5.2} & \textbf{0.759} & \textcolor{successgreen}{\textbf{+0.044}} & 0.710 & 0.965 & 0.831 & 0 \\
DeepSeek-V4-Flash & 0.747 & \textcolor{successgreen}{+0.031} & 0.562 & 0.955 & 0.810 & 0 \\
Qwen3.6-27B & 0.743 & \textcolor{successgreen}{+0.028} & 0.591 & 0.960 & 0.813 & 19 \\
gpt-oss-120b & 0.731 & \textcolor{successgreen}{+0.016} & 0.341 & 0.959 & 0.774 & 2 \\
gpt-oss-20b & 0.724 & \textcolor{successgreen}{+0.009} & 0.455 & 0.956 & 0.797 & 62 \\
Llama-3.3-70B & 0.710 & \textcolor{dangerred}{$-$0.005} & 0.227 & 0.950 & 0.759 & 71 \\
\bottomrule
\end{tabular}
\end{block}
\centering\footnotesize
\textbf{GLM-5.2 thắng rõ:} 0.715 $\rightarrow$ \textbf{0.759}, precision 0.726 $\rightarrow$ \textbf{0.831}, bỏ 71\% báo giả, giữ 96.5\% phát hiện thật (0 lỗi).\\
{\scriptsize Nguồn: kernel Kaggle \texttt{whaleeatu/k8-model-compare}, N=1000, 2026-08-13, 8\,019\,s. Nền 0.715 (subset) khác 0.713 (full-val).}
\end{frame}

% ============================================================
\section{Generation — vòng RAG}
{
\setbeamercolor{background canvas}{bg=white}
\begin{frame}[plain]
    \vfill\centering
    \textcolor{hcmusblue}{\textbf{\Huge Generation — luận giải grounded (vòng RAG)}}
    \vfill
\end{frame}
}

\begin{frame}{Vòng RAG — sinh luận giải grounded}
\centering
\begin{tikzpicture}[node distance=0.45cm and 0.55cm, every node/.append style={inner sep=2.5pt}]
  \node[box, text width=1.6cm] (q) {Tài liệu\\{\tiny query text/.txt}};
  \node[box, text width=1.6cm, right=of q, draw=primaryblue] (ret) {Retrieval\\{\tiny TF-IDF}};
  \node[gbox, text width=2.4cm, right=of ret] (pr) {Prompt (augmented)\\{\tiny template + span}\\{\tiny\textcolor{accentorange}{+ nguồn truy hồi}}};
  \node[obox, text width=1.4cm, right=of pr] (llm) {LLM\\{\tiny GLM-5.2}};
  \node[gbox, text width=1.9cm, right=of llm] (out) {Luận giải\\{\tiny grounded}};
  \node[nbox, text width=1.5cm, below=0.6cm of ret] (corp) {kho nguồn};
  \draw[ar] (q) -- (ret);
  \draw[arO] (ret) -- node[slbl,above]{augment} (pr);
  \draw[ar] (pr) -- (llm);
  \draw[ar] (llm) -- (out);
  \draw[ar] (corp) -- (ret);
\end{tikzpicture}

\vspace{0.1cm}
\begin{columns}[T]
\begin{column}{0.5\textwidth}
\setbeamercolor{block title}{bg=primaryblue,fg=white}
\begin{block}{Vì sao đúng là RAG}
\scriptsize
\begin{itemize}\setlength\itemsep{0pt}
  \item[\ding{51}] \textbf{Query = tài liệu} người dùng nộp (không gõ prompt).
  \item[\ding{51}] \textbf{Retrieve:} TF-IDF lấy đoạn nguồn — bằng chứng.
  \item[\ding{51}] \textbf{Augment:} hệ tự chèn nguồn vào prompt (mắt xích "A").
  \item[\ding{51}] \textbf{Generate:} LLM sinh luận giải neo ngữ cảnh, không bịa.
\end{itemize}
\end{block}
\end{column}
\begin{column}{0.48\textwidth}
\setbeamercolor{block title}{bg=successgreen,fg=white}
\begin{block}{Đo bước sinh (viết lại)}
\scriptsize
\begin{itemize}\setlength\itemsep{0pt}
  \item char-ratio bị-phát-hiện: \textbf{0.176 $\rightarrow$ 0.147}
  \item fidelity (giữ nghĩa): \textbf{0.955}
  \item overlap từ vựng: \textbf{0.139 $\rightarrow$ 0.100}
\end{itemize}
\end{block}
\end{column}
\end{columns}

\vspace{-0.05cm}
\setbeamercolor{block title}{bg=dangerred,fg=white}
\begin{block}{Trung thực về phạm vi}
\scriptsize Đo trên \textbf{2 span / 1 tài liệu} (Gemini, giới hạn quota). Luận giải là \textbf{bổ sung kiến trúc mới}, đã chạy đầu-cuối trong demo (GLM-5.2), chưa full-val. Lõi đo đầy đủ vẫn là retrieval + alignment + verifier.
\end{block}
\end{frame}

% ============================================================
\section{Thước đo}
{
\setbeamercolor{background canvas}{bg=white}
\begin{frame}[plain]
    \vfill\centering
    \textcolor{hcmusblue}{\textbf{\Huge Thước đo}}
    \vfill
\end{frame}
}

\begin{frame}{Thước đo — Retrieval \& Detection}
\small
\begin{columns}[T]
\begin{column}{0.48\textwidth}
\setbeamercolor{block title}{bg=primaryblue,fg=white}
\begin{block}{Recall@k}
\[ R@k = \frac{\#\text{susp có nguồn thật trong top-k}}{\#\text{susp}} \]
\footnotesize \textbf{R@1 = 0.975}: 97.5\% trường hợp nguồn đúng ở hạng 1.
\end{block}
\begin{block}{MRR}
\[ \mathrm{MRR} = \frac{1}{N}\sum \frac{1}{\mathrm{rank}(\text{nguồn thật})} \]
\footnotesize \textbf{0.985} $\rightarrow$ nguồn đúng gần như luôn ở đỉnh.
\end{block}
\end{column}
\begin{column}{0.48\textwidth}
\setbeamercolor{block title}{bg=accentorange,fg=white}
\begin{block}{Precision / Recall (mức ký tự)}
\[ P = \frac{1}{|S|}\sum_{s\in S}\frac{|s\cap(\bigcup R)|}{|s|}, \quad
   R = \frac{1}{|R|}\sum_{r\in R}\frac{|r\cap(\bigcup S)|}{|r|} \]
\footnotesize $R$ = đoạn gold, $S$ = đoạn hệ phát hiện.
\end{block}
\begin{block}{Granularity}
\[ gran = \frac{1}{|R_{det}|}\sum \#\text{đoạn-}S\text{ phủ mỗi đoạn-}R \]
\footnotesize $=1$ lý tưởng; $>1$ = phân mảnh, bị phạt.
\end{block}
\end{column}
\end{columns}
\end{frame}

\begin{frame}{Thước đo — PlagDet \& Verifier}
\small
\begin{columns}[T]
\begin{column}{0.5\textwidth}
\setbeamercolor{block title}{bg=successgreen,fg=white}
\begin{block}{PlagDet (PAN 2015) — thước đo chính}
\[ \mathrm{PlagDet} = \frac{F1}{\log_2(1+gran)}, \; F1 = \frac{2PR}{P+R} \]
\footnotesize Cân bằng P, R rồi chia hình phạt phân mảnh. Hệ đạt \textbf{0.713} $\approx$ \textbf{83\% của trần câu 0.858}.
\end{block}
\end{column}
\begin{column}{0.48\textwidth}
\setbeamercolor{block title}{bg=primaryblue,fg=white}
\begin{block}{Verifier — khử báo giả}
\[ \text{fp\_reduction} = \frac{\#FP\text{ bị bỏ}}{\#FP} \]
\[ \text{tp\_retention} = \frac{\#TP\text{ được giữ}}{\#TP} \]
\footnotesize fp\_reduction cao = bỏ nhiều báo giả; tp\_retention $\approx 1$ (đừng vứt phát hiện thật).
\end{block}
\end{column}
\end{columns}

\vspace{0.15cm}
\setbeamercolor{block title}{bg=softgray,fg=darktext}
\begin{block}{\footnotesize Vì sao tách thước đo?}
\footnotesize Một điểm gộp che điểm yếu. Recall@k cho biết retrieval có đưa đúng nguồn; nếu retrieval trượt thì PlagDet thấp \emph{không phải} lỗi alignment. Tách ra để biết \textbf{khâu nào cần cải thiện}.
\end{block}
\end{frame}

% ============================================================
{
\setbeamercolor{background canvas}{bg=white}
\begin{frame}[plain]
    \vfill\centering
    \textcolor{hcmusblue}{\textbf{\Huge Cảm ơn thầy và các bạn đã lắng nghe}}
    \vfill
\end{frame}
}

\end{document}
